% Purpose: Standalone LaTeX file to draft and preview the Introduction section.
% Function: Includes a minimal preamble and sets the section counter to 0 so it renders as Section 1.

\documentclass[11pt]{article}
\usepackage[utf8]{inputenc}
\usepackage[a4paper,margin=1.0in]{geometry}
\usepackage{amsmath,amssymb,amsthm} 

%\linespread{1.15}
\linespread{2.0} % double space (BMC Medical Research Methodology)

\begin{document}
\setcounter{section}{0} % Forces the following \section to be numbered 1

% --- INTRODUCTION ---
\section{Introduction}
% 1) general background & problem statement
% 2) detail the problem: spillover and spatial dependence
% 3) gap in current literature & motivation
% 4) study contributions and objectives
% 5) paper roadmap

\subsection{Background \& Motivation}

Randomized controlled trials (RCTs) represent the preeminent methodological framework for establishing causal inference when evaluating the efficacy and safety of novel interventions \cite{pocock_clinical_1983, smith_field_2015}. By randomly allocating subjects to intervention or control arms, investigators ensure that measured and unmeasured confounding variables are theoretically balanced in expectation, thereby permitting an unbiased estimation of the average treatment effect. However, in numerous public health and social science contexts, randomizing at the individual level is logistically infeasible or scientifically inappropriate due to the community-level delivery of the intervention. For instance, interventions such as mass media health education campaigns or community-wide vector control mechanisms inherently apply to collective groups rather than discrete individuals. In such scenarios, the cluster randomized trial (CRT) emerges as a rigorous alternative, allocating intact social units or geographically defined regions (e.g., schools, clinics, counties) to study arms \cite{hayes_cluster_2017, donner_design_2000}.

Although CRTs successfully mitigate logistical barriers and attenuate within-cluster treatment contamination, geographically clustered trials introduce complex spatial dynamics that challenge the fundamental statistical assumption of unit independence. Two primary and interrelated spatial phenomena uniquely compromise the internal validity of spatial CRTs: spillover effects and spatial dependence \cite{hayes_cluster_2017, hudgens_toward_2008}. Spillover, or interference, occurs when the application of an intervention in one cluster alters the potential outcomes of subjects residing in neighboring unassigned clusters. This transboundary phenomenon is extensively documented in community-based research, particularly within infectious disease interventions such as the deployment of permethrin-impregnated bed nets \cite{binka_impact_1998, gimnig_effect_2003, hawley_community-wide_2003} and school-based deworming programs \cite{miguel_worms_2004}. Under these conditions, control clusters situated near intervention clusters may receive a protective positive externality, thereby artificially enhancing their observed outcomes. Failing to account for this interference violates the Stable Unit Treatment Value Assumption (SUTVA), frequently precipitating a severe underestimation of the true direct effect of the intervention \cite{anayaizquierdo_spatial_2021}.

Compounding this methodological challenge is spatial dependence, conceptualized by Tobler’s First Law of Geography, which postulates that proximally located entities exhibit greater similarity than distant ones \cite{tobler_computer_1970}. In a spatial CRT, adjacent clusters often share unmeasured environmental exposures or socioeconomic demographics, causing their outcomes to be correlated independent of treatment assignment \cite{cressie_statistics_1993}. Ignoring this underlying spatial structure inherently assumes a greater effective sample size than actually exists, which artificially deflates standard errors, inflates the Type I error rate, and jeopardizes valid statistical inference \cite{campbell_how_2014, ntani_consequences_2021}.

Historically, the biostatistical and epidemiological literature has primarily confronted these challenges through retrospective analytical adjustments rather than via proactive trial design modifications. A systematic review by Jarvis et al. (2017) indicated that spatial location is frequently treated as an afterthought in CRT design \cite{jarvis_spatial_2017}. When researchers do account for spatial phenomena, they typically rely upon spatial variable approaches (such as incorporating distance-based covariates) \cite{hawley_community-wide_2003, lenhart_insecticidetreated_2008, chao_contribution_2015} or spatial modeling techniques that integrate geographic structure into the covariance matrix \cite{alexander_spatial_2003, silcocks_spatial_2010}. While robust, relying solely on statistical correction post-data collection may result in a substantive loss of statistical power and remains highly vulnerable to model misspecification. Consequently, recent methodological literature has begun to advocate for addressing spatial interference proactively at the design stage, demonstrating that strictly spatial sampling strategies like block stratified sampling can minimize estimation variance compared to simple random sampling by physically isolating treatment clusters.

Nevertheless, a critical gap remains regarding the performance of these design strategies under real-world conditions. The current spatial design literature predominantly operates under the assumption of a homogeneous spatial landscape. Empirically, this assumption rarely holds; geographic regions exhibit pronounced, spatially heterogeneous baseline incidence of primary outcomes prior to the introduction of any intervention. Conditions such as Sudden Unexpected Death (SUD) or localized infectious disease outbreaks frequently manifest in distinct geographic hotspots driven by localized socioeconomic determinants \cite{mirzaei_years_2019}. When a spatially structured incidence profile exists, relying on strictly geometric randomization risks inducing severe systematic confounding. An intervention assignment may inadvertently align perfectly with an underlying disease hotspot, inducing structural multicollinearity. This confounds the spatial lag of the treatment with the endogenous spatial lag of the baseline risk, rendering it mathematically intractable for spatial autoregressive models to disentangle the direct intervention effect from the clustered pre-existing burden. Consequently, there is a profound lack of evidence-based guidance detailing how investigators might actively leverage known prior distributions of outcome incidence to optimize spatial treatment allocation.

\subsection{Aims \& Contributions}

The present study addresses this critical methodological gap by shifting the analytical focus from retrospective spatial adjustments toward the rigorous evaluation of incidence-guided treatment sampling designs. By integrating pre-intervention epidemiological data directly into the randomization protocol, researchers can theoretically orthogonalize the treatment assignment from underlying spatial risks, thereby preserving the validity of subsequent causal inferences. The primary objective is to determine the extent to which the incorporation of historical incidence data during the treatment assignment phase minimizes estimation error for the direct intervention effect, particularly in the presence of unobserved spillover and spatial dependence. To achieve this objective, an extensive Monte Carlo simulation study is conducted on a regular spatial lattice, systematically comparing the estimation efficiency of six distinct treatment assignment strategies. This analysis contrasts three strictly spatial sampling methods (block stratification, 2x2 random blocking, and isolation buffer treatment clusters) against three novel covariate-constrained methods that directly utilize baseline incidence data (median incidence grouping, saturation quadrants, and treatment balanced incidence quartiles).

Through this simulation framework, a comprehensive Data Generating Process is formalized utilizing a Spatial Durbin Model (SDM) that explicitly parameterizes both endogenous spatial dependence and exogenous treatment spillovers. The SDM provides a robust theoretical foundation by simultaneously accommodating the direct effects of the intervention, the indirect spillover effects originating from neighboring treated clusters, and the underlying spatial autocorrelation inherent to the outcome variable. Design performance is subsequently evaluated across multiple realistic data generation structures, including independent uniform incidence, spatially correlated incidence via a Simultaneous Autoregressive (SAR) filter, and Poisson count-based rates modeled after North Carolina SUD demographic data. Ultimately, the findings demonstrate that the estimation of intervention effects via rigorous spatial Maximum Likelihood Estimation (MLE) achieves optimal Mean Squared Error (MSE) predominantly when the initial sampling design actively balances underlying baseline risk and ensures sufficient regional variation. By identifying the systematic vulnerabilities of rigid spatial designs in heterogeneous environments, this research offers practical guidance enabling investigators to tailor sampling strategies based on prior incidence to facilitate more efficient and valid causal inference in complex spatial settings.

The remainder of this article is organized as follows. Section 2 delineates the theoretical framework, detailing the spatial structure, the Spatial Durbin Model data generating process, the distinct incidence generation modes, and the specific algorithms underlying the six treatment assignment designs. Section 3 details the simulation study, encompassing the parameter space exploration, evaluation metrics, and comprehensive quantitative results. Section 4 presents an empirical application of these findings to a localized study on Sudden Unexpected Death in Lenoir County, North Carolina. Finally, Section 5 discusses the broader implications of these findings, provides actionable recommendations for researchers designing spatial CRTs, and outlines prevailing limitations alongside requisite avenues for future research.

% --- BIBLIOGRAPHY (from Zotero) ---
% sample citations

% print bibliography
%\bibliographystyle{plain} % alphabetical order
\bibliographystyle{unsrt} % order of appearance
\bibliography{references.bib}

\end{document}